\section{{\it Random Forest}}

{\it Random Forest} merupakan algoritma klasifikasi berbasis metode \textit{ensemble learning} yang menggabungkan hasil dari banyak pohon keputusan (\textit{decision tree}) untuk menghasilkan prediksi yang lebih stabil dan akurat. Algoritma ini dikembangkan oleh Breiman (2001) dengan tujuan meningkatkan performa klasifikasi dibandingkan pohon keputusan tunggal yang rentan terhadap {\it overfitting}. Setiap pohon dalam {\it Random Forest} dilatih menggunakan teknik {\it bootstrap aggregating} atau {\it bagging}, yaitu dengan mengambil sampel acak dari data pelatihan dengan pengembalian, serta memilih subset acak dari fitur saat membangun node pohon. Hasil akhir prediksi diperoleh dengan menggunakan mayoritas suara (\textit{majority voting}) pada klasifikasi atau rata-rata (\textit{mean}) pada regresi (Hastie, Tibshirani, \& Friedman, 2009).
    
    Secara matematis, untuk klasifikasi hasil prediksi dari {\it Random Forest} dapat diformulasikan sebagai berikut.  
    \[
    \hat{y} = \text{mode}\{h_1(x), h_2(x), ..., h_k(x)\}
    \]  
    di mana \( h_i(x) \) adalah hasil prediksi dari pohon ke-\( i \), dan \( k \) adalah jumlah total pohon. Jika mayoritas pohon memprediksi bahwa seorang penyewa termasuk kategori mencurigakan, maka model akan memberikan label mencurigakan. Selain itu, {\it Random Forest} juga dapat menghasilkan nilai probabilitas, yaitu persentase jumlah pohon yang memutuskan kelas tertentu. Sebagai contoh, jika 80 dari 100 pohon memutuskan bahwa penyewa mencurigakan, maka probabilitas penyewa mencurigakan adalah 0.80 atau 80\%.
    
    Sebagai ilustrasi, misalkan terdapat data penyewa dengan fitur lama menginap, metode pembayaran, dan waktu check-in. Data tersebut dimasukkan ke dalam model {\it Random Forest} dengan 100 pohon. Setelah proses prediksi, sebanyak 72 pohon memutuskan penyewa adalah normal dan 28 pohon memutuskan mencurigakan. Maka, probabilitas penyewa diklasifikasikan sebagai mencurigakan sebesar 28\% dan sebagai normal sebesar 72\%. Hasil ini dapat diinterpretasikan bahwa model lebih cenderung mengklasifikasikan penyewa tersebut sebagai normal, meskipun tingkat keyakinannya belum 100\%. Dalam praktik pengawasan, ambang batas (\textit{threshold}) tertentu dapat ditetapkan untuk mengklasifikasikan penyewa sebagai mencurigakan, misalnya jika probabilitas $\geq$60\%, maka akan muncul notifikasi peringatan. 
    
    Keunggulan {\it Random Forest} terletak pada kemampuannya menangani data berdimensi tinggi dan kelas yang tidak seimbang tanpa memerlukan banyak praproses. Selain itu, algoritma ini juga memberikan estimasi pentingnya fitur (\textit{feature importance}) yang bermanfaat untuk interpretasi model dan pengambilan keputusan. Dengan keunggulan tersebut, {\it Random Forest} telah digunakan secara luas dalam deteksi anomali, penipuan, dan klasifikasi risiko dalam berbagai domain, termasuk keamanan siber, kesehatan, dan properti (Liaw \& Wiener, 2002).